



The generic spin Hamiltonian  with the applied magnetic field is
\begin{equation}
\begin{split}
    H   &=\frac{1}{2}\sum_{ ij }  \bm{\sigma}_i \cdot \mathbf{J}_{ij} \cdot \bm{\sigma}_j  - \sum_i \mathbf{h}_i \cdot \bm{\sigma}_i    \; ,\\
         &=\frac{1}{2}\sum_{ ij } \mathbf{J}_{ij}^{\alpha\beta}    \sigma_i^\alpha  \sigma_j^\beta  - \sum_{i\alpha} h_i^\alpha \sigma_i^\alpha    \; ,
\end{split}
\end{equation}
Using the Majorana Parton representation for the spin,
\begin{equation}
    \sigma_i^\alpha = \im c_i^0 c_i^\alpha =  \frac12 \im c_i^{\T} \rmN^\alpha  c_i \; ,
\end{equation},
together with a Lagrange multiplier leads to the Hamiltonian
\begin{equation}
    H   =\frac{1}{2^3}\sum_{ ij }
         \im c_i^{\T} \rmN^\alpha  c_i  \,  \mathbf{J}_{ij}^{\alpha\beta} \,  \im c_j^{\T} \rmN^\beta  c_j
         -\sum_{i\gamma} \Big( h_i^\gamma  \frac{\im}{2} c_i^{\T} \rmN^\gamma  c_i - \lambda_i^\gamma \frac{\im}{4} c_i^{\T} \rmG^\gamma  c_i  \Big)    \; . \label{B:H-full}
\end{equation}
We will consider a general mean-field decoupling for this generic Hamiltonian to show how the Majorana mean-field approximation can be used.

\section{Mean field approximation}

In the mean-field approximation, we decompose the quartic terms in Majorana in bilinear using
\begin{equation}
\begin{split}
        \im c_i^\mu c_i^\nu \,
        \im c_j^ \mup  c_j^\nup
        &\approx \im c_i^\mu c_i^\nu  \rmV_{j}^{ \mup \nup} \!\!\! +   \im  c_j^ \mup   c_j^\nup  \rmV_{i}^{\mu\nu} \!\!\! - \rmV_{i}^{\mu\nu} \rmV_{j}^{ \mup \nup}   , \\
        &- \im c_i^\mu c_j^ \mup    \rmU_{i j }^{\nu\nup} \! - \im  c_i^\nu  c_j^\nup \rmU_{ij}^{\mu \mup }  \! + \rmU_{ij}^{\mu \mup }   \rmU_{i j }^{\nu\nup} , \\
        &+ \im c_i^\mu c_j^\nup   \rmU_{ij}^{\nu \mup } \! +  \im  c_i^\nu  c_j^ \mup  \rmU_{ij}^{\mu\nup} \!   -  \rmU_{ij}^{\mu\nup}  \rmU_{ij}^{\nu \mup } .  \label{eq:MF_decouple}
\end{split}
\end{equation}
with two types of mean-field parameters: $\rmV_i$ the product of two Majorana fermions at the same site, and $\rmU_{ij}$ at different sites, %Which generalizes the expectations value for the polarized and for the Kitaev spin liquid phases
\begin{align}
        \rmV_{i}^{\mu\nu} &= \braket{\im c_i^\mu c_i^\nu } , \quad
        \rmU_{ij}^{\mu\nu} = \braket{\im c_i^\mu c_j^\nu } . \label{eq:MF_parameters}
\end{align}
Where $\rmU^{\mu\nu}$ generalizes the $\mathbb{Z}_2$ gauge field $u_{\braket{ij}^\alpha}=\braket{\im c^\alpha_i c^\alpha_j}$ as the mean-field order parameter for the Kitaev spin liquid phase, and $\rmV_i$ for the polarized phase where $\braket{\sigma^\alpha}\neq 0$.

Then
\begin{align}
\begin{split}
        \im c_i^\mu \rmN^\alpha_{\mu\nu} c_i^\nu \, &
        \im c_j^ \mup \rmN^\beta_{\mup\nup}   c_j^\nup
        \approx \\[8pt]
        &\;\im c_i^\mu \rmN^\alpha_{\mu\nu} c_i^\nu  \rmV_{j}^{ \mup \nup}\rmN^\beta_{\mup\nup}  \!\!\!
        +   \im  c_j^ \mup  \rmN^\beta_{\mup\nup}  c_j^\nup  \rmV_{i}^{\mu\nu}\rmN^\alpha_{\mu\nu}  \!\!\!
        - \rmV_{i}^{\mu\nu}  \rmN^\alpha_{\mu\nu}  \rmV_{j}^{ \mup \nup}  \rmN^\beta_{\mup\nup}  \\[8pt]
        &- \im c_i^\mu \rmN^\alpha_{\mu\nu} \rmU_{i j }^{\nu\nup}  \rmN^\beta_{\mup\nup} c_j^\mup    \!
        - \im  c_i^\nu \rmN^\alpha_{\mu\nu}   \rmU_{ij}^{\mu \mup } \rmN^\beta_{\mup\nup} c_j^\nup \!
        + \rmU_{ij}^{\mu \mup }\rmN^\alpha_{\mu\nu}  \rmU_{i j }^{\nu\nup}  \rmN^\beta_{\mup\nup}
        \\[8pt]
        &+ \im c_i^\mu \rmN^\alpha_{\mu\nu} \rmU_{ij}^{\nu \mup }\rmN^\beta_{\mup\nup}  c_j^\nup    \!
        +  \im  c_i^\nu \rmN^\alpha_{\mu\nu}   \rmU_{ij}^{\mu\nup}\rmN^\beta_{\mup\nup}  c_j^ \mup  \!
        -  \rmU_{ij}^{\mu\nup} \rmN^\alpha_{\mu\nu}  \rmU_{ij}^{\nu \mup }\rmN^\beta_{\mup\nup}
\end{split}\\[6pt]
\begin{split}
        &=
        \im c_i \rmN^\alpha  c_i  \; ( \rmV_{j} , \rmN^\beta )
        + \im c_j \rmN^\beta  c_j  \; ( \rmV_{i} , \rmN^\alpha )
        - ( \rmV_{i} , \rmN^\alpha )\;( \rmV_{j} , \rmN^\beta )   \\[3pt]
        &- \im c_i ( \rmN^\alpha \rmU_{ij} \rmN^{\beta \, \T} ) c_j   \!
        -  \im c_i ( \rmN^{\alpha\, \T} \rmU_{ij} \rmN^{\beta } ) c_j  \!
        + (  \rmU_{ij}  , \rmN^\alpha \rmU_{ij} \rmN^{\beta\, \T} )
        \\[3pt]
        &+ \im c_i ( \rmN^\alpha \rmU_{ij} \rmN^\beta ) c_j   \!
        +  \im c_i ( \rmN^{\alpha\, \T} \rmU_{ij} \rmN^{\beta \, \T} ) c_j  \!
        - (  \rmU_{ij}  , \rmN^\alpha \rmU_{ij} \rmN^\beta )
\end{split}\\[6pt]
\begin{split}
        &=
        \im c_i \rmN^\alpha  c_i  \; ( \rmV_{j} , \rmN^\beta )
        + \im c_j \rmN^\beta  c_j  \; ( \rmV_{i} , \rmN^\alpha )
        + 4 \im c_i ( \rmN^\alpha \rmU_{ij} \rmN^\beta ) c_j  \\[3pt]
        &- ( \rmV_{i} , \rmN^\alpha )\;( \rmV_{j} , \rmN^\beta )
        - 2(  \rmU_{ij}  , \rmN^\alpha \rmU_{ij} \rmN^\beta )  \; .
\end{split}
\end{align}
with matrix inner product $(\rmA , \rmB) = \sum_{\mu\nu} \rmA^{\mu\nu} \rmB^{\mu\nu} = \tr (\rmA^\T \rmB)$ and using that $\rmN^{\alpha\, \T} = - \rmN^\alpha$. Also, from the fully symmetry of $\mathbf{J}_{ij}^{\alpha\beta}$ follows
\begin{equation}
  \frac{1}{2} \sum_{ ij }
         \mathbf{J}_{ij}^{\alpha\beta}    c_i  \; ( \rmV_{j} , \rmN^\beta )
        + \im c_j \rmN^\beta  c_j  \; ( \rmV_{i} , \rmN^\alpha )  =
       \sum_{ ij }
         \mathbf{J}_{ij}^{\alpha\beta}  \im c_i \rmN^\alpha  c_i  \; ( \rmV_{j} , \rmN^\beta ) \; .
\end{equation}
Applying this decoupling into the Hamiltonian resumes into equation \eqref{eq:HMF}.

\begin{comment}

To compute the mean-field parameters auto-consistently one needs to diagonalize the mean-field Hamiltonian. We do it by the unitary transformation $\mathbb{U}$ where we first pair up the Majorana fermions into complex fermions and then unitary diagonalize the complex fermionic Hamiltonian,
\begin{equation}
    c = \sqrt{2} \mathbb{U} \, \left( \begin{array}{c}  d^{\phantom{\dagger}} \\ d^\dagger \end{array} \right)  , \quad
    \mathbb{U} = \Big( \begin{array}{cc} \mathbb{U}_{<} & \mathbb{U}_{>}
    \end{array}\Big)\; .
\end{equation}
Where the columns of $\mathbb{U}_{<(>)}$ are the negative (positive) eigenvectors of the mean-field Hamiltonian. The mean-field ground state is the state annihilated by all $d_\lambda$ operators, from which follows
\begin{equation}
\begin{split}
        \braket{ \,  \im   c_\I \,  c_\J  \, }
        &=\im  \sum_{\varepsilon_\lambda \geq 0 } \big( \mathbb{U}^{\dagger} \big)_{\lambda,\I} \, \big( \mathbb{U}\, \big)_{\J,\lambda}
         = \im \big(\mathbb{U}_{>}^\pd  \mathbb{U}_{>}^\dagger  \big)_{\J\I}
        \\&=\im  \sum_{\varepsilon_\lambda \leq 0 } \big( \mathbb{U}\, \big)_{\I,\lambda} \, \big( \mathbb{U}^{\dagger} \big)_{\lambda,\J} \,
         = \im \big(\mathbb{U}_{<}^{\pd}  \mathbb{U}_{<}^\dagger  \big)_{\I\J}
         \; , \label{eq:MF_eq}
\end{split}
\end{equation}
where we used the multi-index for site and fermionic flavor $\I=(i\mu)$, $\J=(j\nu)$. Restricting $j$ to be either equal or nearest-neighbor to $i$ we find the mean-field parameters in .


The energy expectation value
\begin{equation}
\begin{split}
    \braket{H_{\mathrm{MF}}  } &= %\rmC
    \frac{1}{8}\!\sum_{ij}  \mathbf{J}_{ij}^{\alpha\beta} \Big[ \tr\!\left( \rmV_{i}^\T \rmN^\alpha \right)\tr\!\left( \rmV_{j}^\T \rmN^\beta \right)+2\tr\!\left( \rmU_{ij}^\T \rmN^\alpha   \rmU_{ij} \rmN^\beta \right)   \Big] \\
    &+\frac{1}{4}\!\sum_{i, \gamma}  \!\Big[  -  2  h_i^\gamma  \tr\!\left( \rmV_{i}^\T \rmN^\gamma \right)  + \lambda_i^\gamma \tr\!\left( \rmV_{i}^\T \rmG^\gamma  \right) \Big]  ,
\end{split}
\end{equation}

\end{comment}
\begin{comment}

It is convenient to define
\begin{equation}
    \tilde{\lambda}_i^\gamma = \lambda_i^\gamma + h_i^\gamma
\end{equation}
so that
\begin{equation}
\begin{split}
    \braket{H_{\mathrm{MF}}  } &= %\rmC
    \frac{1}{8}\!\sum_{ij}  \mathbf{J}_{ij}^{\alpha\beta} \Big[ \tr\!\left( \rmV_{i}^\T \rmN^\alpha \right)\tr\!\left( \rmV_{j}^\T \rmN^\beta \right)+2\tr\!\left( \rmU_{ij}^\T \rmN^\alpha   \rmU_{ij} \rmN^\beta \right)   \Big] \\
    &+\frac{1}{4}\!\sum_{i, \gamma}  \!\Big[  -    h_i^\gamma  \tr\!\left( \rmV_{i}^\T \rmN^\gamma \right)  + \tilde{\lambda}_i^\gamma \tr\!\left( \rmV_{i}^\T \rmG^\gamma  \right) \Big]  ,
\end{split}
\end{equation}

Since the matrices $\rmV$  are $2\im \mathbb{1}_4$ plus an antissimetric matrix them can be expanded as
\begin{equation}
    \rmV_{j}(\lambda) = 2\im \mathbb{1}_4 + \sum_\gamma w^\gamma_j(\lambda)  \rmN^\gamma + \sum_\gamma \tilde{w}^\gamma_j(\lambda) \rmG^\gamma
\end{equation}
where the projection onto physical states implies that  $\tilde{w}^\gamma_j\vert_{\lambda_c} = 0$.

Using the inner product  $\left(\rmA,\rmB\right):=\tr\!\left( \rmA^\T \rmB \right) $ for the spin, isospin  and $\rmN=(\rmN-\rmG)/2 $ matrices
\begin{align}
    \left( \rmN^\alpha , \rmN^\beta \right) &= 4 \delta^{\alpha\beta}, \quad
    \left( \rmG^\alpha , \rmG^\beta \right) = 4 \delta^{\alpha\beta}, \quad
    \left( \rmN^\alpha , \rmG^\beta \right) = 0 \; , \\
    \left( \rmN^\alpha , \rmN^\beta \right) &= 2 \delta^{\alpha\beta}, \quad
    \left( \rmN^\alpha , \rmG^\beta \right) = -2 \delta^{\alpha\beta} .
\end{align}
We have
\begin{equation}
    \left( \rmV_{j} , \rmN^\gamma  \right) = 4 w^\gamma_j, \quad
    \left( \rmV_{j} , \rmN^\gamma  \right) = 2\big( w^\gamma_j- \tilde{w}^\gamma_j\big) , \quad
    \left( \rmV_{j} , \rmG^\gamma  \right) = 4 \tilde{w}^\gamma_j .
\end{equation}
Then
\begin{equation}
\begin{split}
    \braket{H_{\mathrm{MF}}  } &= %\rmC
    \!\sum_{ij}  \mathbf{J}_{ij}^{\alpha\beta} \Big[ \frac{1}{2}\big( w^\alpha_i- \tilde{w}^\alpha_i\big)\big( w^\beta_j- \tilde{w}^\beta_j\big)+\frac{1}{4}\tr\!\left( \rmU_{ij}^\T \rmN^\alpha   \rmU_{ij} \rmN^\beta \right)   \Big] \\
    &+\!\sum_{i, \gamma}  \!\Big[  -  h_i^\gamma  w^\gamma_i + \tilde{\lambda}_i^\gamma \tilde{w}^\gamma_i \Big]  ,
\end{split}
\end{equation}

Idea: take the expansion in power of $\tilde{w}^\gamma_i \sim \delta\lambda = (\lambda - \lambda_c)$. The first order term contains $\lambda_c$. Equivalently, take the derivative with
$\delta\lambda$


$\;$

\vspace{15mm}


Note that
\begin{equation}
    \frac{\partial \braket{H_{\mathrm{MF}}  }}{\partial \rmV_{i}} =
    \frac{1}{4}\sum_\gamma \Big[   \mathbf{J}_{\braket{ij}_\gamma }^{\alpha\beta}
    \rmN^\alpha \tr\!\left( \rmV_{j}^\T \rmN^\beta \right) -    2 h_i^\gamma  \rmN^\gamma   + \lambda_i^\gamma  \rmG^\gamma   \Big]
\end{equation}
we find projecting on $\rmG^\gamma$ that
\begin{align}
    -2\sum_{\beta}\mathbf{J}_{\braket{ij}_\gamma }^{\gamma \beta}
      \tr\!\left( \rmV_{j}^\T \rmN^\beta \right) +   4 h_i^\gamma     + 4 \lambda_i^\gamma   &= 0
\end{align}
and writing $(\mathbb{J}_i)^{\gamma\beta}=\mathbf{J}_{\braket{ij}_\gamma }^{\gamma\beta}$ ,
\begin{equation}
    \mathbb{J}=\left(
\begin{array}{ccc}
 J_x+K_x & 0 & 0 \\
 0 & J_y+K_y & 0 \\
 0 & 0 & J_z+K_z \\
\end{array}
\right),
\end{equation}
makes clear that only the term with $\beta=\gamma$ is non-zero (at least for $\Gamma^\prime=0$) and then
\begin{equation}
    \lambda_i^\gamma = -  h_i^\gamma + \frac{1}{2} \mathbf{J}_{\braket{ij}_\gamma }^{\gamma \gamma}
     \tr\!\left( \rmV_{j}^\T \rmN^\gamma \right) \; ,
\end{equation}
where $j$ is such that it is in a $\braket{ij}_\gamma$ bond with respect to $i$.


\end{comment}

\clearpage
\section{Fourier transform}

Assuming a translational invariant mean-field ground state, we can simplify the mean-field Hamiltonian and the self-consistent equation.
Let us first rewrite the Hamiltonian as a sum over unit cells
\begin{equation}
    H_{\mathrm{MF}}=
    \frac{\im}{4} \sum_{r,\gamma}   c_{r,A}^{\T}  2 \rmA_{rA,r_\gamma B}  c_{r_\gamma,B} +
    \frac{\im}{4}  \sum_{r}\Big[ c_{r,A}^{\T}\rmB_{r,A}c_{r,A}+c_{r,B}^{\T}\rmB_{r,B}c_{r,B}\Big] -
    \rmC ,
\end{equation}
where the matrices in the Hamiltonian are
\begin{align}
\begin{split}
    \rmC&=\frac{1}{16}\sum_{r\gamma }   \mathbf{J}_{rA,r_{\gamma} B}^{\alpha\beta} \Big[
     \tr\!\left( \rmV_{rA}^\T \rmN^\alpha \right)\tr\!\left( \rmV_{r_\gamma B}^\T \rmN^\beta \right)
    + 2 \,  \tr\!\left( \rmU_{rA,r_\gamma B}^\T \rmN^\alpha   \rmU_{rA,r_\gamma B} \rmN^\beta \right)   \Big]
    \; ,
\end{split}    \\
\begin{split}
    \rmB_{rA}&= \sum_{\gamma}\!\Big[  \frac{1}{4}  \mathbf{J}_{rA,r_{\gamma} B}^{\alpha\beta}
    \rmN^\alpha \tr\!\left( \rmV_{r_{\gamma} B}^\T \rmN^\beta \right) -    h_{rA}^\gamma  \rmN^\gamma  + \lambda_{rA}^\gamma  \rmG^\gamma \Big]
    \; ,
\end{split} \\
\begin{split}
    \rmB_{rB}&= \sum_{\gamma}\!\Big[  \frac{1}{4} \mathbf{J}_{rB,r_{-\gamma} A}^{\alpha\beta}
    \rmN^\alpha \tr\!\left( \rmV_{r_{-\gamma} A}^\T \rmN^\beta \right) -    h_{rB}^\gamma  \rmN^\gamma  + \lambda_{rB}^\gamma  \rmG^\gamma \Big]
    \; ,
\end{split} \\
\rmA_{rA,r_\gamma B} &= \frac{1}{2} \mathbf{J}_{rA,r_{\gamma} B}^{\alpha\beta}   \rmN^\alpha \rmU_{rA,r_{\gamma} B} \rmN^\beta \; .
\end{align}
The Fourier transformed modes are
\begin{align}
    c_{\mathbf{k},\sigma}^\mu &= \frac{1}{\sqrt{2N}} \sum_{r} \; c_{r,\sigma}^\mu e^{-\im \mathbf{k}\cdot\mathbf{R}_r} \; , \\
    c_{r,\sigma}^\mu &= \sqrt{\frac{2}{N}} \sum_{\mathbf{k}}^{\prime}  \; \Big[
    c_{\mathbf{k},\sigma}^{\mu\pd\!} e^{\im \mathbf{k}\cdot\mathbf{R}_r}+c_{\mathbf{k},\sigma}^{\mu\dagger} e^{-\im \mathbf{k}\cdot\mathbf{R}_r}
    \Big]\; .
\end{align}
Under the translation invariance assumption, the Hamiltonian reduces to
\begin{align}
    \!\!H_{\mathrm{MF}}
    &=\!\!
    \frac{1}{2}\sum_{\mathbf{k}} \!\!  c_{\mathbf{k},A}^{\dagger}  \sum_{\gamma} \Big[2\im \rmA_{\gamma} e^{\im \mathbf{k}\cdot \bm{n}_\gamma} \Big]  c_{\mathbf{k},B}  +\frac{1}{2}
    \!\sum_{\mathbf{k}}\im\Big[ c_{\mathbf{k},A}^{\dagger}\rmB_{A} c_{\mathbf{k},A}^\pd
    +c_{\mathbf{k},B}^{\dagger}  \rmB_{B} c_{\mathbf{k},B}^\pd \Big] -
    \rmC ,  \\
    &=-      \rmC  + \frac{1}{2}
    \sum_{\mathbf{k}}
    \left(\begin{array}{c} c_{\mathbf{k},A} \\ c_{\mathbf{k},B}
    \end{array}\right)^{\dagger}  \!
    \left(\begin{array}{cc}
         \im \rmB_{A} & \im \sum_{\gamma} \rmA_{\gamma} e^{\im \mathbf{k}\cdot \bm{n}_\gamma}  \\
         -\im \sum_{\gamma} \rmA_{\gamma}^{\T} e^{-\im \mathbf{k}\cdot \bm{n}_\gamma} & \im \rmB_{B}
    \end{array}\right)\!
    \left(\begin{array}{c} c_{\mathbf{k},A} \\ c_{\mathbf{k},B}
    \end{array}\right) \\
    &=-      \rmC  + \frac{1}{2} \!
    \sum_{\mathbf{k}}
    \left(\begin{array}{c} \bar{d}_{\mathbf{k}} \\ d_{\mathbf{k}}
    \end{array}\right)^{\T}  \!
    \left(\begin{array}{cc}
         \text{diag}\big(\varepsilon_n(\mathbf{k}) \big) & 0\\
         0& \text{diag}\big(-\varepsilon_n(\mathbf{k}) \big)
    \end{array}\right)\!
    \left(\begin{array}{c} d_{\mathbf{k}} \\ \bar{d}_{\mathbf{k}}
    \end{array}\right)\\
    &= -\rmC + \sum_{n, \mathbf{k}} \varepsilon_n(\mathbf{k}) d_{\mathbf{k}}^\dagger d_{\mathbf{k}}^\pd -\frac{1}{2}\!   \sum_{n, \mathbf{k}} \varepsilon_n(\mathbf{k}) \; .
\end{align}
where
\begin{align}
\begin{split}
    \rmC&= \frac{N}{16} \sum_{\gamma }   \mathbf{J}_{\gamma}^{\alpha\beta} \Big[
      \tr\!\left( \rmV_{A}^\T \rmN^\alpha \right)\tr\!\left( \rmV_{B}^\T \rmN^\beta \right)
    + 2 \,  \tr\!\left( \rmU_{\gamma}^\T \rmN^\alpha   \rmU_{\gamma} \rmN^\beta \right)   \Big]
    \; ,
\end{split}    \\
\begin{split}
    \rmB_{A}&= \sum_{\gamma}\!\Big[ \frac{1}{4} \mathbf{J}_{\gamma}^{\alpha\beta}
    \rmN^\alpha \tr\!\left( \rmV_{B}^\T \rmN^\beta \right) -    h_{A}^\gamma  \rmN^\gamma  + \lambda_{A}^\gamma  \rmG^\gamma \Big]
    \; ,
\end{split} \\
\begin{split}
    \rmB_{B}&= \sum_{\gamma}\!\Big[  \frac{1}{4}  \mathbf{J}_{\gamma}^{\alpha\beta}
    \rmN^\alpha \tr\!\left( \rmV_{A}^\T \rmN^\beta \right) -    h_{B}^\gamma  \rmN^\gamma  + \lambda_{B}^\gamma  \rmG^\gamma \Big]
    \; ,
\end{split} \\
\rmA_{\gamma} &= \frac{1}{2}\mathbf{J}_{\gamma}^{\alpha\beta}   \rmN^\alpha \rmU_{\gamma} \rmN^\beta
\end{align}

We can define complex fermions such that
\begin{align}
    c_{\mathbf{k}}^\pd &= \mathbb{U}_{<}^\pd(\mathbf{k}) \, d_\mathbf{k}^\pd +  \mathbb{U}_{>}^\pd(\mathbf{k}) \, \bar{d}_\mathbf{k}\; , \\[4pt]
    c_{\mathbf{k}}^\dagger &=  d_\mathbf{k}^\T \mathbb{U}_{>}^\dagger(\mathbf{k}) +   d_\mathbf{k}^\dagger \mathbb{U}_{<}^\dagger(\mathbf{k}) \; ,
    \end{align}
using overline on $d$ to represent the column vector\footnote{We reserve $d_{\mb k}^\dagger$ for the row vector obtained from Hermitean conjugation of the column vector $d_{\mb k}$.} with entries $(\bar{d}_{\mb k})_\rho =(d_{\mb k}^\dagger)_\rho $.
The unitary transformation is split with half of the columns in $\mathbb{U}_{<}$ and the other half in $\mathbb{U}_{>}$, such that $\mathbb{U}_{<} = \left( \mathbb{U}_{<} \; \mathbb{U}_{>} \right)$.
These matrices are related by charge conjugation symmetry due to the Majorana fermion setup which satisfies $c_{-\mathbf{k}}= (c_{\mathbf{k}}^\dagger)^\T = \mathbb{U}_{>}^\ast(\mathbf{k}) \, d_\mathbf{k}^\pd +  \mathbb{U}_{<}^\ast(\mathbf{k}) \, \bar{d}_\mathbf{k} $, forcing us to define $\mathbb{U}_{</>}(-\mathbf{k})=\mathbb{U}_{>/<}^\ast(\mathbf{k})$.
The unitary transformation diagonalizes the Hamiltonian
    \begin{align}
    \mathbb{U}^\dagger
    \left(\begin{array}{cc}
         \im \rmB_{A} & \im \sum_{\gamma} \rmA_{\gamma} e^{\im \mathbf{k}\cdot \bm{n}_\gamma}\!\!\!  \\
         \!\!-\im \sum_{\gamma} \!\rmA_{\gamma}^{\T} e^{-\im \mathbf{k}\cdot \bm{n}_\gamma} & \im \rmB_{B}
    \end{array}\right)\! \mathbb{U} &=
    \left(\begin{array}{cc}
         \bm{\varepsilon} & 0\\
         0& -\bm{\varepsilon}
    \end{array}\right) ,
\end{align}
where $\bm{\varepsilon}$ is the $8\times 8$ diagonal matrix $\text{diag}\big(\varepsilon_n(\mathbf{k}) \big)$.

Using the unitary transformations $\mathbb{U}(\mb{k})$, we can find the mean-field parameters (for unit cells $\mathbf{R}'=\mathbf{R}+\mathbf{n}_\gamma$)
\begin{align}
    \braket{ \im c_{r} c_{r_\gamma }^{T} }
    &= \frac{1}{N}\sum_r \braket{ \im c_{r} c_{r_\gamma }^{T} } \\
    &= \frac{2}{N^2}\sum_{r,\mathbf{k},\mathbf{k}'} \braket{ \im c_{\mathbf{k}}^\pd  c_{\mathbf{k}' }^\dagger } e^{\im \mathbf{R}_{r}\cdot(\mathbf{k}-\mathbf{k}')} e^{-\im \mathbf{k}'\cdot \mathbf{n}_\gamma}  \\
    &= \frac{2}{N}\sum_{\mathbf{k} } \braket{ \im c_{\mathbf{k}}^\pd  c_{\mathbf{k} }^\dagger } e^{-\im \mathbf{k} \cdot \mathbf{n}_\gamma}  \\
    &= \frac{2}{N}\sum_{\mathbf{k} } \im \mathbb{U}_{<}^\pd(\mathbf{k}) \mathbb{U}_{<}^\dagger(\mathbf{k})  e^{-\im \mathbf{k} \cdot \mathbf{n}_\gamma}   \; ,
    %\\    &= \frac{1}{N}\sum_{\mathbf{k}  } \im \Big(     \mathbb{U}_{>}^\ast(\mathbf{k}) \mathbb{U}_{>}^\T(\mathbf{k})  e^{\im \mathbf{k} \cdot \mathbf{n}_\gamma}  +    \mathbb{U}_{<}^\pd(\mathbf{k}) \mathbb{U}_{<}^\dagger(\mathbf{k})  e^{-\im \mathbf{k} \cdot \mathbf{n}_\gamma}      \Big) .
\end{align}
where $\braket{ \im c_{r} c_{r_\gamma }^{T} }$ is an $8\times8$ matrix with entries $\braket{ \im c_{r,\sigma}^\mu  c_{r_\gamma,\sigma' }^\nu}$.
We use for Majorana fermions $c_r^T=c_r^\dagger$.
This expression allows us to compute the mean-field parameters in the translational invariant case.
The algorithm to solve the self-consistent equation is the same as described in the main text, sec.~\ref{sec:4-algorithm}, for the problem in the real lattice.

As is evidenced by the number of mean-field parameters, the problem with translational symmetry is much simpler and less computationally costly than the one in the real lattice.
This allows us (in more reasonable time) to test the limits of convergence for many more parameters $J, K, \Gamma, \bm{h}$, system size $Lx,Ly$, and ways to enforce the constraints with the Lagrange multiplier $\lambda$. Some of this benchmark is shown in the appendix \ref{App-translational}.










\section{Lagrange Multiplier}

The mean-field Hamiltonian is
\begin{align}
    H_{\mathrm{MF}} &=\frac{\im}{4} \sum_{ij}  c_i^{\T}   \rmA_{ij}  c_j + \frac{\im}{4} \sum_i  c_i^{\T}   \rmB_{i}  c_i- \rmC \; , \\
    \rmC &= \frac{1}{8} \sum_{ij}  \mathbf{J}_{ij}^{\alpha\beta} \Big[
       \braket{ \rmV_{i} \, , \,  \rmN^\alpha } \braket{ \rmV_{j} \, , \,  \rmN^\beta }
       +2\braket{ \rmU_{ij} \, , \,  \rmN^\alpha   \rmU_{ij} \rmN^\beta }   \Big]   \; ,  \label{eq:C} \\
    \rmB_{i}&=\! \sum_{j}  \mathbf{J}_{ij}^{\alpha\beta}
    \rmN^\alpha \braket{ \rmV_{j} \, , \,  \rmN^\beta }
    -\sum_{\gamma}\!\Big( 2h_i^\gamma  \rmN^\gamma  - \lambda_i^\gamma  \rmG^\gamma \Big)
    \; , \label{eq:B} \\
    \rmA_{ij} \!&= 2 \mathbf{J}_{ij}^{\alpha\beta}  \rmN^\alpha \rmU_{ij} \rmN^\beta     \; ,  \label{eq:A}
\end{align}
where  $\rmN^\gamma=(\rmN^\gamma-\rmG^\gamma) $ and $\gamma = 1,2,3$ for $\{\rmN^\gamma\}$ and $\{\rmG^\gamma\}$ being two linear independent real representations of the $\mathfrak{su}(2)$ algebra inside the fundamental representation of $\mathfrak{so}(4)=\mathfrak{su}(2)\times\mathfrak{su}(2)$. Thus, the span of $\{ \rmN^\gamma, \rmG^\gamma \}$ is the $6$-dimensional set of $4\times 4$ antissimetric matrices.
The matrix representation of $\rmN^\gamma, \rmG^\gamma $ is given in appendix \ref{App-Parton}, see eq.~\eqref{A:MG}.
These matrices are orthogonal with respect to the inner product $\braket{\rmA \, , \,  \rmB} :=\tr\!\left( \rmA^\T \rmB \right) $ :
\begin{align}
    \braket{\rmN^\alpha , \rmN^\beta} &=  \delta^{\alpha\beta}, \quad
    \braket{ \rmG^\alpha , \rmG^\beta } =  \delta^{\alpha\beta}, \quad
    \braket{ \rmN^\alpha , \rmG^\beta } = 0 \; , \\
    \braket{ \rmN^\alpha , \rmN^\beta } &=  \delta^{\alpha\beta}, \quad
    \braket{ \rmN^\alpha , \rmG^\beta } = - \delta^{\alpha\beta} .
\end{align}

Since the matrices $\rmV$  are $2\im \mathbb{1}_4$ plus an antissimetric matrix they can be expanded as
\begin{equation}
    \rmV_{j}(\lambda) = 2\im \mathbb{1}_4 + \sum_\gamma \mathrm{v}^\gamma_j(\lambda)  \rmN^\gamma + \sum_\gamma \bar{\mathrm{v}}^\gamma_j(\lambda) \rmG^\gamma\;,
\end{equation}
with coefficients given by the inner-products
\begin{equation}
    \braket{ \rmV_{j} , \rmN^\gamma } =  \mathrm{v}^\gamma_j, \quad
    \braket{ \rmV_{j} , \rmN^\gamma } = \big( \mathrm{v}^\gamma_j - \bar{\mathrm{v}}^\gamma_j\big) , \quad
    \braket{ \rmV_{j} , \rmG^\gamma } =  \bar{\mathrm{v}}^\gamma_j .
\end{equation}
Using these expression, the energy functional is   %\reminder{check factors of 2 and 4}
\begin{align}
\begin{split}
    \braket{H_{\mathrm{MF}}  } &= %\rmC
    \frac{1}{8}\!\sum_{ij}  \mathbf{J}_{ij}^{\alpha\beta} \Big[
    \braket{ \rmV_{i} \, , \,  \rmN^\alpha } \braket{ \rmV_{j} \, , \,  \rmN^\beta }
       +2\braket{ \rmU_{ij} \, , \,  \rmN^\alpha   \rmU_{ij} \rmN^\beta }    \Big] \\ &
    +\frac{1}{4}\!\sum_{i, \gamma}  \!\Big[  -  2  h_i^\gamma \braket{ \rmV_{i} \, , \,  \rmN^\gamma }  + \lambda_i^\gamma \braket{ \rmV_{i} \, , \,  \rmG^\gamma } \Big]  ,
\end{split}\\
\begin{split}
     &=
   \frac{1}{8}\!\sum_{ij}  \mathbf{J}_{ij}^{\alpha\beta} \Big[  \big( \mathrm{v}^\alpha_i- \bar{\mathrm{v}}^\alpha_i\big)\big( \mathrm{v}^\beta_j- \bar{\mathrm{v}}^\beta_j\big)+2\braket{ \rmU_{ij} \, , \,  \rmN^\alpha   \rmU_{ij} \rmN^\beta }     \Big] \\
    &+\frac{1}{4}\!\sum_{i, \alpha}  \!\Big[ -2h_i^\alpha  \big( \mathrm{v}^\alpha_i - \bar{\mathrm{v}}^\alpha_i \big) + \lambda_i^\alpha \bar{\mathrm{v}}^\alpha_i \Big]   \; ,
\end{split}\\
\begin{split}
     &=
   \frac{1}{4}\!\sum_{ij}\mathbf{J}_{ij}^{\alpha\beta}  \braket{ \rmU_{ij} \, , \,  \rmN^\alpha   \rmU_{ij} \rmN^\beta }     +\frac{1}{4}\!\sum_{i, \alpha}  \!\Big[ -2{h}_i^\alpha  \big( \mathrm{v}^\alpha_i - \bar{\mathrm{v}}^\alpha_i \big) + \lambda_i^\alpha \bar{\mathrm{v}}^\alpha_i \Big]  ,  \\
   &+
   \frac{1}{8}\!\sum_{i}  \big( \mathrm{v}^\alpha_i- \bar{\mathrm{v}}^\alpha_i\big) \sum_{j}  \mathbf{J}_{ij}^{\alpha\beta}\big( \mathrm{v}^\beta_j- \bar{\mathrm{v}}^\beta_j\big).
\end{split}
\end{align}
The matrix element $\mathbf{J}_{ij}^{\alpha\beta}$ is only non-zero for $j$ in a bond with $i$, so it restricts to the three terms in the sum over $\gamma$ in $\braket{ij}_\gamma$, as  $i$ can be in either sublattice here (to not confusing with) the notation where $\braket{ij}$ forces $i/j\in A/B$  I will write this sum as a sum of nearest neighbors $\sum_{j\in V_i}$.
From this, we can identify the effective magnetic field
\begin{align}
        \braket{H_{\mathrm{MF}}  }
     &=
   \frac{1}{4}\!\sum_{ij}\mathbf{J}_{ij}^{\alpha\beta}  \braket{ \rmU_{ij} \, , \,  \rmN^\alpha   \rmU_{ij} \rmN^\beta }     +\frac{1}{4}\!\sum_{i, \alpha}  \!\Big[ -2\tilde{h}_i^\alpha  \big( \mathrm{v}^\alpha_i - \bar{\mathrm{v}}^\alpha_i \big) + \lambda_i^\alpha \bar{\mathrm{v}}^\alpha_i \Big]  ,  \\
   \tilde{h}_i^\alpha &= h_i^\alpha - \frac{1}{2}
   \!\sum_{j \in V_i} \! J^{\alpha\beta}_{ij}  \big( \mathrm{v}^\beta_j- \bar{\mathrm{v}}^\beta_j\big)
   \; . \label{B:effective-field}
\end{align}



\section{Partially polarized phase}
We can solve the mean-field equation in the phase where all generalized Kitaev type parameters are zero $U_{ij}^{\alpha\beta}=0$.
Let us consider first the uniform case
\begin{equation}
     \mathrm{v}^\gamma_i=\mathrm{v}^\gamma \; , \quad
     \bar{\mathrm{v}}^\gamma_i=\bar{\mathrm{v}}^\gamma \;.
\end{equation}
The energy density $F = \braket{H_{\mathrm{MF}}  }/2N$  is
\begin{align}
\begin{split}
    F     &=
   \!\sum_{\gamma}  \Big[ \; \mathbf{J}_{\gamma}^{\alpha\beta}\frac{1}{2}\big( \mathrm{v}^\alpha-\bar{\mathrm{v}}^\alpha\big)\big( \mathrm{v}^\beta-\bar{\mathrm{v}}^\beta\big)  -  h^\gamma  \big( \mathrm{v}^\gamma - \bar{\mathrm{v}}^\gamma \big)  + \lambda^\gamma \bar{\mathrm{v}}^\gamma \Big]  ,
\end{split}
\end{align}
where $\mathbf{J}_{\gamma}^{\alpha\beta} \equiv  \mathbf{J}_{\braket{ij}_\gamma}^{\alpha\beta} $ and $\sum_\gamma \equiv \sum_{j \in V_i }$ for the set $V_i=\{i+\delta_x,i+\delta_y,i+\delta_z\}$ of N.N. of $i$.
Instead of directly imposing the minimization of $F$ w.r.t to $\lambda$ that enforces $\bar{\mathrm{v}}=0$, let us vary $F$ w.r.t. $\bar{\mathrm{v}}$ and find the critical value for the Lagrange multiplier
\begin{align}
\begin{split}
    \frac{\partial F}{\partial \bar{\mathrm{v}}^\beta}      &=
   \! -\Big(\sum_{\gamma}  \mathbf{J}_{\gamma}^{\alpha\beta} \Big)\big( \mathrm{v}^\alpha-\bar{\mathrm{v}}^\alpha\big) +  h^\beta    + \lambda^\beta   \Rightarrow \boxed{\;\lambda^\alpha = -h^\alpha + \mathbb{J}^{\alpha\beta}\big( \mathrm{v}^\beta-\bar{\mathrm{v}}^\beta\big)  \, } \, ,
\end{split}
\end{align}
where
\begin{equation}
    \mathbb{J}^{\alpha\beta} = \frac{\partial^2 F}{\partial \bar{\mathrm{v}}^\alpha \partial \bar{\mathrm{v}}^\beta} =
    \sum_{\gamma}  \mathbf{J}_{\gamma}^{\alpha\beta} =
    \left(\begin{array}{ccc}
       3J + K& \Gamma & \Gamma \\
       \Gamma & 3J + K  & \Gamma \\
       \Gamma &\Gamma &3J + K
    \end{array}\right)^{\alpha\beta} \; .
\end{equation}
Note that $\det \left(\mathbb{J}^{\alpha\beta}\right) = (-\Gamma +3 J+K)^2 (2 \Gamma +3 J+K)$ is negative close to the pure FM Kitaev limit $K<0$ so the energy is a minimum.

This minimization procedure implies that the Lagrange multiplier is equal to $\lambda = h-2 \tilde{h}$.
Plugging this value for the Lagrange multiplier in the energy density
\begin{align}
\begin{split}
    F     &=
    \frac{1}{2} \mathbb{J}^{\alpha\beta}  \big( \mathrm{v}^\alpha-\bar{\mathrm{v}}^\alpha\big)\big( \mathrm{v}^\beta-\bar{\mathrm{v}}^\beta\big)  -  h^\alpha   \mathrm{v}^\alpha +\mathbb{J}^{\alpha\beta}\bar{\mathrm{v}}^\alpha \big( \mathrm{v}^\beta-\bar{\mathrm{v}}^\beta\big) ,
\end{split}\\
\begin{split}
         &=
    \frac{1}{2} \mathbb{J}^{\alpha\beta}  \big( \mathrm{v}^\alpha\mathrm{v}^\beta-\bar{\mathrm{v}}^\alpha\bar{\mathrm{v}}^\beta\big)   -  h^\alpha   \mathrm{v}^\alpha  .
\end{split}
\end{align}
From that, we immediately see that if $\det \left(\mathbb{J}^{\alpha\beta}\right)<0$  the energy is minimized for $\bar{\mathrm{v}}=0$. Minimizing the energy w.r.t. $\mathrm{v}^\alpha$ we find
\begin{align}
\begin{split}
    \frac{\partial F}{\partial \mathrm{v}^\alpha}      &=
   \mathbb{J}^{\alpha\beta}  \mathrm{v}^\beta
   -  h^\alpha
   \Rightarrow
   \boxed{\;\mathrm{v}^\alpha = \big(\mathbb{J}^{-1}\big)^{\alpha\beta} \; h^\beta   \, } \, ,
\end{split}
\end{align}
and the classical polarized energy is
$F=-\frac{1}{2} \big(\mathbb{J}^{-1}\big)^{\alpha\beta} \; h^\alpha h^\beta$. Note however that this solution breaks down in the region where $\det \left(\mathbb{J}^{\alpha\beta}\right)>0$, i.e. $\Gamma > -\frac{K}{2}-\frac{3J}{2}$, which gives a explanation for the transition at $\Gamma/(-K)=0.5$ in the $h\Gamma$ phase  diagram for $J=0$.

In the inhomogeneous case, the Lagrange multiplier is
\begin{equation}
    \lambda_i^\alpha = - h^\alpha_i + \frac{1}{2} \!\! \sum_{j \in V_i} J^{\alpha\beta}_{ij}
    \big( \mathrm{v}^\beta_j-\bar{\mathrm{v}}^\beta_j\big)  \; ,
\end{equation}
which gives the free energy
\begin{align}
\begin{split}
    2N \, F
    %&=      \!\sum_{i} \Big[  -  h_i^\gamma    \mathrm{v}^\gamma_i \!+\!\! \sum_{j(\gamma) \in V_i}\!\!J^{\alpha\beta}_{\braket{ij}_\gamma} \frac{1}{2}\Big(        \big( \mathrm{v}^\alpha_i- \bar{\mathrm{v}}^\alpha_i\big)\big( \mathrm{v}^\beta_j- \bar{\mathrm{v}}^\beta_j\big)     +     \bar{\mathrm{v}}^\alpha_i\big( \mathrm{v}^\beta_j- \bar{\mathrm{v}}^\beta_j\big)    \Big)\Big], \end{split}\\
%\begin{split}    &=  \!\sum_{i} \Big[  -  h_i^\gamma    \mathrm{v}^\gamma_i \!+\frac{1}{2}\!\! \sum_{j(\gamma) \in V_i}\!\!J^{\alpha\beta}_{\braket{ij}_\gamma} \big( \mathrm{v}^\alpha_i+    \bar{\mathrm{v}}^\alpha_i\big) \big( \mathrm{v}^\beta_j- \bar{\mathrm{v}}^\beta_j\big)  \Big],\end{split} \\ \begin{split}
   &=  \!-\sum_{i}  h_i^\gamma    \mathrm{v}^\gamma_i \!+\frac{1}{2} \sum_{ij}J^{\alpha\beta}_{\braket{ij}} \big( \mathrm{v}^\alpha_i \mathrm{v}^\beta_j-
    \bar{\mathrm{v}}^\alpha_i  \bar{\mathrm{v}}^\beta_j\big)  \,  .
\end{split}
\end{align}




\section{Kitaev spin liquid phase}\label{sec:B-lag}
Treating the mean-field parameters $\rmU, \rmV$ as independent variables, the minimization procedure for $\rmV$ gives the same expression for the Lagrange multiplier. Allowing a correction, I will write the Lagrange multiplier as
\begin{equation}
    \lambda_i^\alpha = - h^\alpha_i + \frac{1}{2}  \sum_{j \in V_i} J^{\alpha\beta}_{ij}
    \big( \mathrm{v}^\beta_j-\bar{\mathrm{v}}^\beta_j\big) + \delta \lambda_i  \; , \label{B:lagrange-multiplier}
\end{equation}
where $\delta h_i$ is such that the solution of the self-consistent equation satisfies the constraint $\braket{ \rmV_{i} \, , \,  \rmG^\gamma } =0$. As there is no way to solve the self-consistent equations in the general case, then one cannot compute $ \delta \lambda_i$.   A not elegant solution for this is: for each fixed value of $\{\delta \lambda_i\}$ one test if it solves the self-consistent equations and computes the value of $\braket{ \rmV_{i} \, , \,  \rmG^\gamma }$, iterating this process one can find the solution by sampling a set of values one.

In the homogeneous case, one can solve it by brute force. Define functions $\alpha_1, \alpha_2, \alpha_3$
\begin{equation}
    \delta \bm{\lambda} = \alpha_1 \mathbf{a}  + \alpha_2 \mathbf{b} + \alpha_3 \mathbf{c} \; ,
\end{equation}
These are functions of the field and the couplings.
From trial and error, I found that $\delta \bm{\lambda}$ is small and scales linearly with $\mathbf{h}$. The problem reduces to find $\alpha$ such that
\begin{equation}
    \delta \bm{\lambda} = \alpha \mathbf{h} \;
\end{equation}
implies in the constraint fulfillment
\begin{equation}
    \braket{\im c_i^T \rmG^\alpha  c_i }
=
\tr ( \rmV_i^T \rmG^\alpha)
= \frac{
\rmV_i^{0\alpha}-\rmV_i^{\beta \gamma } }{2}
    =\frac{
    \braket{\im c_i^0 c^\alpha_i }  - \braket{ \im c^\beta_i c^\gamma_i}}{2} = 0 \; .
\end{equation}
We need to find $\bm{\lambda} = -\mb{h}_{\text{eff}} + \delta \bm{\lambda}$ such that the mean-field solution satisfies the these constraints.

To measure the violation from $\tr ( \rmV_i^T \rmG^\alpha)=0$ we define
\begin{equation}
    \Delta \rmV_i(\alpha) := \frac{1}{\sqrt{3} }||~\bar{\mathrm{v}}_i~||
    = \sqrt{\;  \frac{1}{3}\sum_{\gamma = 1}^{3} \;\braket{\rmV_i\,,\,\rmG^\gamma }^2 \; }
\end{equation}
Where the inner product is $\braket{\rmV_i\,,\,\rmG^\gamma }=\tr ( \rmV_i^T \rmG^\alpha)$.
This violation is zero states that are physical as satisfies $\tr ( \rmV_i^T \rmG^\alpha)=0$.
The normalization is chosen so that the maximal value is one, with the factor $1/3$ due to three terms in the $\gamma$ sum.  %, normalization of $\rmG: \braket{ \rmV, \rmG^\gamma } =  \bar{\mathrm{v}}^\gamma$).



\begin{figure}
    \centering    \includegraphics[width=.7\textwidth]{Chapters/Appendix/Figures/mf/cv-lambda.pdf}
    %\includegraphics[width=.7\textwidth]{Chapters/Appendix/Figures/mf/en-lambda.pdf}
    \caption{The optimal value for the Lagrange multiplier is $\delta \lambda = -0.11 \mathbf{h}$. }
    \label{fig-delta-v}
\end{figure}



The constraint violation for %$\rmV = 2 \im  \mathbb{1}_4 + \sum_\gamma \mathrm{v}^\gamma  \rmN^\gamma + \bar{\mathrm{v}}^\gamma \rmG^\gamma $ is
$\rmV^{\mu\nu}=\left<ic^\mu c^{\nu}\right>$ is
\begin{equation}
 (\Delta \rmV)^2 %:= \frac{1}{3} \sum_\gamma (\bar{\mathrm{v}}^\gamma)^2
 =\dfrac{
 \left(\dfrac{\left<ic^0c^{{x}}\right>+\left<ic^{{y}}c^{{z}}\right>}{2}\right)^2+
 \left(\dfrac{\left<ic^0c^{{y}}\right>+\left<ic^{{z}}c^{{x}}\right>}{2}\right)^2+
 \left(\dfrac{\left<ic^0c^{{z}}\right>+\left<ic^{{x}}c^{{y}}\right>}{2}\right)^2
 }{3\; }\;.
\end{equation}
In the isotropic case, this quantity reduces to the expected expression $\Delta \rmV = \frac{1}{2}\left|\left<ic^0c^{{\gamma}}\right>+\left<ic^{{\alpha}}c^{{\beta}}\right>\right|$ with normalization $ \frac{1}{2}$.
%From numerical evidence, I found the maximal value for $|\bar{\mathrm{v}}_i|$ to be one so that the constraint violation is bound to


The bound for the constraint violation is
\begin{align*}
    &0\leq || \bar{\mathrm{v}}_i||^2
\leq
 || {\mathrm{v}}_i||^2 + || \bar{\mathrm{v}}_i||^2
=
\frac{1}{4^2}\sum_\gamma  \big(
\braket{\rmV_i\,,\,\rmN^\gamma }^2 +\braket{\rmV_i\,,\,\rmG^\gamma }^2 \big)
\\& =
\frac{1}{2} \Big(
\braket{\im c^0 c^1 }^2+
\braket{\im c^0 c^2 }^2+
\braket{\im c^0 c^3 }^2+
\braket{\im c^1 c^2 }^2+
\braket{\im c^2 c^3 }^2+
\braket{\im c^3 c^1 }^2
\Big)
\\& \leq
\frac{1}{2} \Big( \!
\braket{(\im c^0 c^1)^2 }+
\braket{(\im c^0 c^2)^2 }+
\braket{(\im c^0 c^3)^2 }+
\braket{(\im c^1 c^2)^2 }+
\braket{(\im c^2 c^3)^2 }+
\braket{(\im c^3 c^1)^2 }\!
\Big)
\\&
%=\frac{1}{2} \Big( \braket{ \bm{\sigma}^2 } + \braket{ \bm{\tau}^2 }\Big)
= 3 \;   \\[3pt]
& \Rightarrow \quad 0 \leq \Delta \rmV \leq 1 \; .
\end{align*}
%where the spin $\frac{1}{2}\bm{\sigma}$ and the pseudo spin $\frac{1}{2}\bm{\tau}$ squared are $\frac{3}{2}$
%8 (w[2, 1]^2 + w[3, 1]^2 + w[3, 2]^2 + w[4, 1]^2 + w[4, 2]^2 +    w[4, 3]^2)
%The same equation holds for the magnetization

Now let us turn to analyze how much the constraint is violated for some values of $\delta \lambda$ or $\alpha$.
Numerically $\alpha=0$ already gives a small constraint violation in the KSL region
\begin{equation}
    \alpha = 0 \quad \Rightarrow \quad  \Delta \rmV_i(\alpha)  \lesssim 10^{-2} \; .
\end{equation}
For a more precise solution, one can try different values for $\alpha$ such that $\Delta V$ is minimized.
Here I will consider the homogeneous case with the same $\alpha$ in both sublattices.
For example, for  $\Gamma=0.1$  one can take  $ \alpha \approx -0.11$ to get $ \Delta \rmV < 10^{-5}$ for any value of $h$ as shown in Fig. \ref{fig-delta-v}.

The optimum value for the Lagrange multiplier  is found by fixing all the parameters of the model and varying the Lagrange multiplier indexed by $\alpha$. For each value of the Lagrange multiplier, we independently solve the mean-field self-consistent equation and found the corresponding  constraint violation $\Delta V$.  The value for $\alpha$ is determined a posteriori by finding which minimizes the constraint $\Delta V$. \textit{This solution is \textit{not} elegant}, since we need many trials to determine the value for the Lagrange multiplier in contrast to a self-consistent solution also for the Lagrange multiplier (we tried, but did not found such solution that produced tigher contraints). Fortunetly, for the translation invariant configuration, by this procedure give a almost constant value for the parameter $\alpha$. As we change the value of $\Gamma$, the mean-field parameter $\alpha$ can be well approximated by the affine function  $\alpha(\Gamma)=-0.2+1.13 \Gamma$.  In Fig.~\ref{fig-CV-Gamma}, we show the value of the constraint violation $\Delta V(\alpha)$ for some values of $\Gamma$ and shown that it is minimal for the value $\alpha = \alpha(\Gamma)$. This approximation is surprisingly good, and it is enough for the parameter region that we are interested in.

\begin{figure}
    \centering\hspace*{-10mm}\includegraphics[width=.7\textwidth]{Chapters/Appendix/Figures/mf/lag-m-h.01_.3.pdf}\\
    \includegraphics[width=.8\textwidth]{Chapters/Appendix/Figures/mf/constraint-violation-Gamma-rescaled.pdf}
    \caption{(a) $\alpha(\Gamma)$ for two values of the magnetic field. The dependence in the magnetic field is very weak and in the KSL phase a good approximation is $\alpha(\Gamma)=-0.2 + 1.1 \, \Gamma $.
    (b) Constraint violation vs Lagrange multiplier for different values of $\Gamma$.
    Numerical fit gives linear dependence in $\Gamma$ for the linear coefficient  $\alpha(\Gamma)=-0.2+1.13 \Gamma$.}
    \label{fig-CV-Gamma}
\end{figure}




%\pagebreak

%\section{Phase diagram}

%\begin{figure}
%    \centering    \includegraphics{sections/figures_new/Fig3_Wp.pdf}
%    \centering    \includegraphics{sections/figures_new/Magnetization_out_of_plane.pdf}
%    \centering    \includegraphics{sections/figures_new/DeltaW.pdf}
%    \centering    \includegraphics{sections/figures_new/DeltaW_Real.pdf}
%\end{figure}













